\documentclass[border=10pt]{standalone}
\usepackage{circuitikz}

\begin{document}
\begin{circuitikz}[scale=2, transform shape, line width=1pt, american voltages]
\color{white}

    % Op-amp node
    \node[op amp] (op) at (3,0) {};

    % Non-inverting (+) input connected to ground
    \draw (op.+) node[ground] {};

    % Input branch: Ground -> Vin -> R1 -> Node X -> Inverting (-) input
    \draw (0,-1.5) node[ground] {}
          to[vsource, l=$V_{in}$] ++(0, 2)
          to[R, l=$R_1$] ++(1.8, 0) node[circ] (X) {} node[above left] {$X$}
          -- (op.-);

    % Output wire with a junction dot for feedback connection
    \draw (op.out) -- ++(0.4,0) node[circ] (out_node) {}
          to[short, -o] ++(0.5,0) node[right] {$V_{out}$};

    % Feedback path aligned using orthogonal coordinates (X |- ...) and (out_node |- ...)
    \draw (X) -- (X |- 0, 1.5)
          to[C, l=$C_1$] (out_node |- 0, 1.5)
          -- (out_node);

\end{circuitikz}
\end{document}
